\section{Effective capacities: a successful class and its boundary}
\label{net:effective}
These results analyze an earlier continuous-random-cut rule.  The general
hybrid in Theorem~\ref{prod:upper} does not require its structural incentive
argument.  All networks below are finite directed acyclic $s$--$t$ graphs,
every edge lies on an $s$--$t$ path, and each edge has one distinct fixed owner.
The topology and owners are fixed during every report comparison.  Parallel
edges are permitted unless a graph is explicitly called simple.
Write $F(c)$ for maximum flow, $T=F(c)$, and $v_c(S)=F(c\ind_S)$.
Let $\mathcal C$ be the inclusion-minimal structural edge cutsets and
$K=|\mathcal C|$; this family is independent of capacities, including zero
capacities.  Then $F(c)=\min_{C\in\mathcal C}\sum_{e\in C}c_e$.

Define $D_e(c)$ as the greatest flow through edge $e$ among all feasible
$s$--$t$ flows, without requiring their total value to be maximal.
For independent uniform $U_e\in[0,1]$, choose the first cut $W$ minimizing
$\sum_{e\in C}D_e(c)U_e$, using a fixed priority on $\mathcal C$, and put
\begin{equation}\label{net:base-rule}
 z_e(c)=D_e(c)\Prob(e\in W).
\end{equation}
These are deterministic expectations; the uniforms are an integration device.
A useful graph formula, also proving continuity, is
$D_e=\min\{c_e,F_{s,\mathrm{tail}(e)},F_{\mathrm{head}(e),t}\}$,
omitting a prefix or suffix with coincident terminals.  An upstream path and
a downstream path cannot share a vertex, since that would create a cycle
through $e$.  Scaled maximum prefix and suffix flows therefore join through
$e$ and attain the formula; the reverse bounds are immediate.

\begin{lemma}[Effective-capacity preservation and budget]
\label{net:preservation}
On every such DAG, $D\le c$, each $D_e$ is nondecreasing in every report,
$0\le D_e\le T$, and replacing $c$ by $D(c)$ preserves the complete feasible
flow set of every coalition.  Consequently \eqref{net:base-rule} is
nonnegative and coalition-covering.  For $K\ge2$ let $R_K>e$ solve
$R_K\ln(R_K/e)=\ln K$ and set
\[
 U(K)=R_K+1/\ln R_K.
\]
Then $\sum_ez_e(c)\le U(K)T$; for $K=1$ the rule is efficient.
Here $U(K)\sim\ln K/\ln\ln K$.  No incentive assertion is made in this lemma.
\end{lemma}
\begin{proof}
Inclusion of feasible flow sets proves monotonicity of $D$.  A DAG has no
circulations, so path decomposition bounds edge flow by total flow and hence
by $T$.  Every coalition flow is grand-network feasible, so it respects all
$D_e$; the reverse inclusion uses $D\le c$.
For every selected cut and every coalition,
$\sum_{e\in S\cap W}D_e\ge F(D\ind_S)=v_c(S)$.
The equality follows from preservation and max-flow/min-cut on the fixed
graph after giving nonmembers capacity zero.  Taking expectations proves
coverage.  If $T=0$, all $D_e$ and payments are zero.

For $T>0$ put $a_e=D_e/T\in[0,1]$ and $\rho_C=\sum_{e\in C}a_e$.
Some cut $C_0$ has $\rho_{C_0}=1$.  Its random score is at most one, so the
winning cut has $X_W=\sum_{e\in W}a_eU_e\le1$.
For $t>0$, concavity of $x\mapsto x^{a_e}$ gives
\[
 \E e^{-tX_C}\le\prod_{e\in C}(\E e^{-tU_e})^{a_e}
 \le t^{-\rho_C},\qquad
 \Prob(X_C\le1)\le e^t t^{-\rho_C}.
\]
Set $t=\rho_C$.  Since $(e/s)^s$ decreases for $s>1$, a union bound gives
$\Prob(\rho_W>s)\le\min\{1,K(e/s)^s\}$ for $s\ge e$.
Beyond $R_K$ the logarithm of the second bound has derivative
$-\ln s\le-\ln R_K$.  Therefore
\[
 \frac{\sum_ez_e}{T}=\E\rho_W
 \le R_K+\int_{R_K}^\infty e^{-(s-R_K)\ln R_K}\,ds=U(K).
\]
Cut scores may overlap; no independence between cuts was used.
With one cut, its $D$-sum is $T$.  Finally the defining equation implies
$R_K\sim\ln K/\ln\ln K$, and $1/\ln R_K=o(R_K)$.
\end{proof}

\subsection{Series-parallel networks and co-cut independence}
A two-terminal series-parallel (TTSP) network is built from directed edge
leaves by series and parallel composition of components with disjoint edges
and internal vertices.  Fix its decomposition tree.  Component flow satisfies
$F_{A\parallel B}=F_A+F_B$ and $F_{A;B}=\min(F_A,F_B)$.
Minimal cuts form disjoint alternatives at a series node and unions of one
cut from each child at a parallel node.  Thus
$K_{A;B}=K_A+K_B$ and $K_{A\parallel B}=K_AK_B$.

\begin{theorem}[TTSP incentives]\label{net:ttsp}
On every fixed TTSP network, the rule \eqref{net:base-rule} is own-report
monotone at every nonnegative real profile, and hence satisfies coverage and
the $U(K)T$ budget bound with binding-underreport incentives.
\end{theorem}
\begin{proof}
For an edge $e$, let $L_e$ be the minimum of $F_B(c)$ over the sibling
components $B$ at all series ancestors of $e$, with empty minimum $+\infty$.
Then
\begin{equation}\label{net:ttsp-cap}
 D_e=\min(c_e,L_e).
\end{equation}
For attainability, route exactly this amount through the leaf.  At a parallel
ancestor use zero flow in its other child; at a series ancestor route the same
amount through the sibling by scaling its maximum flow.  Each ancestor's
bound is respected.  Conversely, every path carrying flow through $e$ crosses
each such series sibling, giving all the displayed upper bounds.

If distinct edges $e,h$ share a minimal cut, their lowest common ancestor
must be parallel: at a series node a minimal cut lies wholly in one child.
Formula~\eqref{net:ttsp-cap} shows that $D_h$ depends on $c_e$ only if $e$
lies in a sibling at a series ancestor of $h$.  This is impossible for a pair
sharing a cut.  Thus $D_h$ is independent of $c_e$ for co-cut pairs.

Fix $c_{-e},U_{-e}$ and write $x=c_e$, $d(x)=D_e(c)$.  Put
\[
 A=\min_{C\ni e}\sum_{h\in C\setminus\{e\}}D_hU_h,\qquad
 B(x)=\min_{C\not\ni e}\sum_{h\in C}D_hU_h,
\]
with $B=+\infty$ if no excluding cut exists.
Co-cut independence makes $A$ constant in $x$; $B$ and $d$ are nondecreasing.
Integration over the single uniform $U_e$ gives
\begin{equation}\label{net:conditional}
 D_e\Prob(e\in W\mid U_{-e})=\min\{d(x),(B(x)-A)_+\}.
\end{equation}
When $d>0$ a tie between the containing and excluding scores occurs at at
most one value of $U_e$, hence has measure zero.  The formula also holds for
$d=0$ and $B=+\infty$ under their evident conventions.
Its right side is nondecreasing in $x$.  Average over $U_{-e}$.
\end{proof}
The same proof applies to any fixed cut representation with a normalization
that preserves every coalition, is coordinatewise nondecreasing and bounded
by $T$, and has co-cut independence.  The condition is sufficient, not necessary.
Serial layers of parallel edges recover multi-glove games.  Their lower
bound and Theorem~\ref{net:ttsp} give sharp worst-case TTSP order
$(1+o(1))\ln K/\ln\ln K$.  This lower construction is inherited; the general
upper and simple-DAG lower are now Theorems~\ref{prod:upper} and~\ref{dag:sharp}.

\begin{proposition}[A global cap does not replace effective capacities]
\label{net:global-cap}
Globally replacing each edge by $\min(c_e,T)$ before applying the continuous
random-cut rule can violate own-report monotonicity even on TTSP graphs
with three structural cuts and arbitrarily many owners.
\end{proposition}
\begin{proof}
Use the flow function
$F(x,a,b,d_1,\ldots,d_n)=\min\{x+\min(a,b),\sum_jd_j\}$.
For $n\ge4$, set $\epsilon=1/10$, $a=\epsilon$, $b=n+\epsilon$,
$d_j=1$, and compare $x_0=n-2\epsilon$ with $x_1=n-\epsilon$.
Then $T_0=n-\epsilon$, $T_1=n$.  Under the global cap only $b$ is changed,
to $T$.  If $S_n$ is a sum of $n$ independent unit uniforms and
$A_T=\min(\epsilon U,TV)$ is independent, the focal payment is
\[
 z(x,T)=\E\min\{x,(S_n-A_T)_+\}.
\]
Write $\delta=n-T\in\{\epsilon,0\}$, $t=\epsilon+\delta$, so $x=n-t$.
For $0\le u\le1$, simplex integration gives
$\E(u-S_n)_+=u^{n+1}/(n+1)!$; also $n-S_n$ has the same law as $S_n$.
Since $0\le A_T\le\epsilon\le t<1$, subtraction of the two positive parts
therefore yields
\[
 z(x,T)=\frac n2-\E A_T+
 \frac{\E A_T^{n+1}-\E(t-A_T)^{n+1}}{(n+1)!}.
\]
The survival law $(1-u/\epsilon)(1-u/T)$ on $[0,\epsilon]$ gives
\[
 \E A_T^r=\frac{\epsilon^r}{r+1}
 -\frac{r\epsilon^{r+1}}{(r+1)(r+2)T},\quad r\ge1.
\]
Thus the mean increases by $\epsilon^3/[6n(n-\epsilon)]$ between profiles.
Bounding the two positive remainder terms separately gives
\[
 z(x_1,T_1)-z(x_0,T_0)
 \le-\frac{\epsilon^3}{6n(n-\epsilon)}
       +\frac{2(2\epsilon)^{n+1}}{(n+1)!}<0.
\]
For the last inequality, the ratio of the positive bound to
$\epsilon^3/(6n^2)$ is
$r_n=6n^2 2^{n+2}\epsilon^{n-2}/(n+1)!$.
It equals $64/125<1$ at $n=4$ and decreases, since
$r_{n+1}/r_n=2\epsilon(n+1)^2/[n^2(n+2)]<1$.
Effective capping instead fixes the detour capacity at $a$, as required by
Theorem~\ref{net:ttsp}.  This is a failure of a specific formula, not a
lower bound for all mechanisms.
\end{proof}

\subsection{A bundled Wheatstone bridge works, but composition fails}
Let $H_k$ have vertices $s,u,v,q$ and edges
$x:s\to u$, $y:s\to v$, $z_j:u\to q$ ($1\le j\le k$),
$w:v\to q$, and $h:u\to v$.  These are distinct fixed owners.
Its four minimal cuts are
$\{x,y\}$, $\{z_1,\ldots,z_k,w\}$,
$\{y,z_1,\ldots,z_k,h\}$, and $\{x,w\}$.

\begin{proposition}[Bundled bridge]\label{net:bundle}
For every $k\ge1$, effective-capacity payments are own-report monotone on
$H_k$, at every nonnegative real profile, and pay at most $U(4)T$.
\end{proposition}
\begin{proof}
Writing $C=\sum_jz_j$, the effective capacities are
\[
 \begin{aligned}
 d_x&=\min(x,C+\min(h,w)),&d_y&=\min(y,w),\\
 d_{z_j}&=\min(z_j,x),&d_w&=\min(w,y+\min(x,h)),\\
 d_h&=\min(h,x,w).
 \end{aligned}
\]
Each formula is the maximum amount that can be routed through its specified
edge: the displayed upstream/downstream bounds are necessary and attained
by routing that amount along the indicated direct and bridge paths.
Put $X=d_xU_x$, $Y=d_yU_y$, $Z=\sum_jd_{z_j}U_{z_j}$,
$W=d_wU_w$, $H=d_hU_h$.  Formula~\eqref{net:conditional}, applied to the
four cuts and integrating only the focal uniform, gives the integrands.
Here the containing intercept may depend on the report: only the conditional
probability identity is used, and monotonicity is checked directly below.
The resulting expressions are
\[
 \begin{aligned}
 q_x&=\min(d_x,Z+\min((W-Y)_+,H)),\\
 q_w&=\min(d_w,Y+\min((X-Z)_+,H)),\\
 q_y&=\min(d_y,(W-\min((X-Z)_+,H))_+),\\
 q_{z_j}&=\min(d_{z_j},(X-Z_{-j}-\min((W-Y)_+,H))_+),\\
 q_h&=\min(d_h,(X-Z)_+,(W-Y)_+).
 \end{aligned}
\]
Here $Z_{-j}$ excludes the focal summand.  For example, for $z_j$ the
least containing score without $z_j$ is $Z_{-j}+\min(W,Y+H)$ and the least
excluding score is $X+\min(Y,W)$, which gives the stated cancellation.
The other expressions follow by the same subtraction of the four cut scores.
When $x$ grows, only $d_x,Z,W,H$ can increase and $Y$ stays fixed; all occur
positively in $q_x$.  When $w$ grows, $d_w,Y,X,H$ can increase and $Z$ is
fixed.  When $y$ grows, only $d_y,W$ can increase in its expression.
When $z_j$ grows, only $d_{z_j},X$ can increase; $Z_{-j},W,Y,H$ are fixed.
When $h$ grows, $d_h,X,W$ can increase and $Y,Z$ remain fixed.
Every integrand is therefore own-report monotone.  Average and apply
Lemma~\ref{net:preservation}.
\end{proof}
This is an unbounded-owner, four-cut class outside TTSP, not a theorem about
arbitrarily many interacting bridge structures.

\begin{theorem}[Actual expected-payment failures]
\label{net:wheatstone-failure}
There are fixed edge-owned DAGs with five structural cuts and arbitrarily
many owners on which \eqref{net:base-rule} decreases an owner's expected
payment when only its own report increases.  There are also simple DAGs
with growing vertex and owner counts with the same property.  In both
constructions the two grand values are one and all capacities are positive.
\end{theorem}
\begin{proof}
Add an edge $p:s\to q$ in parallel with $H_k$, then two parallel edges
$r_1,r_2:q\to t$.  The structural cuts are the four bridge cuts with $p$
adjoined, and $\{r_1,r_2\}$.  Set
\[
 \epsilon=1/20,\quad x_0=9/10,\quad x_1=19/20,\quad y=\epsilon,
 \quad z_j=w=h=p=2,\quad r_1=r_2=1/2.
\]
Only $x$ changes.  Both values are one.  For $x\in[x_0,x_1]$ the effective
capacities are
\[
 D_x=x,\ D_y=\epsilon,\ D_{z_j}=D_h=x,\ D_w=x+\epsilon,
 \ D_p=1,\ D_{r_i}=1/2.
\]
The incoming $x$ and $x+y$ cuts and downstream capacity one give the upper
bounds; routes through an individual $z_j$, through $h,w$, and through $p$
attain them.  Let independent uniforms define
\[
 Y=\epsilon U_y,\quad W=(x+\epsilon)U_w,\quad H=xU_h,\quad
 Z=x\sum_{j=1}^kU_{z_j},\quad P=U_p,\quad
 R=(U_{r_1}+U_{r_2})/2,\quad A=\min(Y,W).
\]
The containing-cut score without $x$ is $P+A$; the least excluding score
is $\min\{P+Z+\min(W,Y+H),R\}$.  Thus its actual payment is $\E q_k(x)$,
where
\begin{equation}\label{net:wheatstone-q}
 q_k(x)=\min\{x,[\min\{Z+\min(H,(W-Y)_+),R-P-A\}]_+\}.
\end{equation}
Compare $I(x)=\E\min\{x,(R-P-A)_+\}$.  The integrands agree when $Z\ge1$
and are otherwise between zero and one, so
\begin{equation}\label{net:ideal-error}
 I(x)-\Prob(Z<1)\le z_x(x)\le I(x),\qquad
 \Prob(Z<1)\le\frac1{k!x^k}.
\end{equation}
The second inequality bounds the relevant portion of the unit cube by the
nonnegative simplex of coordinate sum below $1/x$; it is not an equality.

For $b=x+\epsilon$ write
\[
 M_0(b)=1,\qquad
 M_j(b)=\E A^j=\frac{\epsilon^j}{j+1}
       -\frac{j\epsilon^{j+1}}{(j+1)(j+2)b},\quad j\ge1,
 \qquad P_4(t,b)=\sum_{j=0}^4\binom4j(-1)^jt^{4-j}M_j(b).
\]
For $0\le a\le\epsilon$, convolution of independent uniforms of widths
$1/2,1/2,1$ (or simplex inclusion-exclusion) gives
\[
 \E(R-P-a)_+=\frac{(1-a)^4-2(1/2-a)^4}{6}.
\]
At $x_0,x_1$, its truncated part is
$\E(R-P-a-x)_+=(1-x-a)^4/6$.
Consequently
\[
 I(x)=\frac{P_4(1,b)-2P_4(1/2,b)-P_4(1-x,b)}6,
 \qquad I(x_0)-I(x_1)=\frac{553}{136800000}.
\]
By \eqref{net:ideal-error}, the actual payment drop is at least
$553/136800000-(10/9)^k/k!$.
At $k=10$ this is
$97726081746439/30050502681578400000>0$; the error decreases thereafter.
This proves the first assertion for every $k\ge10$.

For the simple version, replace each $z_j$ by a private two-edge path with
both capacities two and distinct owners.  Replace each $r_i$ by a private
two-edge path with both capacities $1/2$.  The other edges are unchanged.
There are $k+7$ vertices and $2k+9$ owners; there are no parallel edges.
The effective capacities of each pair equal the original effective capacity,
and \eqref{net:wheatstone-q} holds with
\[
 Z=x\sum_{j=1}^k\min(U_j,V_j),\qquad
 R=\bigl(\min(U'_1,V'_1)+\min(U'_2,V'_2)\bigr)/2.
\]
Each minimum has density $2(1-u)\le2$, so the error in
\eqref{net:ideal-error} is now at most $2^k/(k!x^k)$.
Each downstream path has density $4-8r$ on $[0,1/2]$.
Convolution gives
\[
 F_R(r)=8r^2-\tfrac{32}{3}r^3+\tfrac83r^4\quad(0\le r\le1/2),
 \qquad 1-F_R(r)=\tfrac83(1-r)^4\quad(1/2\le r\le1).
\]
Since $P$ is uniform, $g(a)=\E(R-P-a)_+=\E(R-a)_+^2/2$.
Integrating $g''=1-F_R$, with $g(0)=5/72$, $g'(0)=-1/3$, yields
\[
 g_0(a)=\tfrac5{72}-\tfrac a3+\tfrac{a^2}2-\tfrac23a^4
               +\tfrac8{15}a^5-\tfrac4{45}a^6\quad(0\le a\le1/2),
 \qquad g_1(a)=\tfrac4{45}(1-a)^6\quad(1/2\le a\le1).
\]
Using $A\le1/20$ and $A+x\ge9/10$ therefore gives
\[
 I(x)=\E g_0(A)-\tfrac4{45}\E(1-x-A)^6,
\]
where every moment is the rational $M_j(b)$ above.  Expansion gives
\[
 I(x_0)=\frac{11809546127}{191520000000},\quad
 I(x_1)=\frac{14797313}{240000000},\quad
 I(x_0)-I(x_1)=\frac{1290353}{191520000000}.
\]
The actual drop is at least
$1290353/191520000000-(20/9)^k/k!$.
At $k=14$ this equals
$700533736635095043154291/118415105665157830969440000000>0$,
and it increases thereafter.  The structural cut count is
$2^{k+1}+6$: each of the two cut types involving the $z$ bundle chooses one
edge of each path, the other two upstream types are unchanged, and the
final pair contributes four cuts.  This proves the simple-DAG assertion.
\end{proof}
The obstruction concerns an actual integrated payment, not only a conditional
integrand.  Coverage and the cut-count budget still hold by
Lemma~\ref{net:preservation}; incentives fail.  In particular, the successful
bundled bridge rule is not preserved by arbitrary series/parallel composition
with other modules.  This does not preclude other mechanisms on these graphs.

\section{The own-report envelope: a linear upper bound and growing obstruction}
\label{net:envelope}
For the rule \eqref{net:base-rule}, its actual own-report monotone envelope is
\[
 y_e(c)=\sup_{0\le t\le c_e}z_e(t,c_{-e}).
\]
It is automatically own-report monotone and dominates the covering vector $z$.
The supremum is attained and defines a continuous function: the graph formula
above makes every $D_e$ continuous.  The conditional formula
$z_e=\E\min\{D_e,(B-A)_+\}$ makes $z_e$ continuous by dominated
convergence, with $A,B$ the containing/excluding scores at the given profile.
Writing $t=\lambda c_e$ with $0\le\lambda\le1$ gives a maximum of a
continuous function over a fixed compact interval, which is continuous.
These regularity observations are not needed for the budget proof.
The next bounds concern this specific repair; they are not bounds on the
optimal subsidy among all truthful mechanisms.

\begin{lemma}[A shared-uniform comparison]\label{net:uniform-comparison}
For finite nonnegative coefficient vectors $a,b$ and independent unit
uniforms, if $A=\sum_i a_i$, $B=\sum_i b_i$, then
\[
 A\Prob\Bigl(\sum_i a_iU_i\le\sum_i b_iU_i\Bigr)\le B.
\]
\end{lemma}
\begin{proof}
If $A\le B$, the assertion is immediate.  Otherwise put
$N=\sum_i(b_i-a_i)_+$ and $W=\sum_i|a_i-b_i|=A-B+2N>0$.
Reflect uniforms whose coefficient $a_i-b_i$ is negative.  The event becomes
$\sum_i|a_i-b_i|V_i\le N$, where $N<W/2$.
The resulting positive weighted uniform sum has a symmetric unimodal density
about $W/2$.  To check this directly, a centered uniform is even and
nonincreasing on the positive half-line.  If $f$ has these properties,
convolution with the centered uniform on $[-r,r]$ has derivative
$[f(x+r)-f(x-r)]/(2r)\le0$ for $x\ge0$ almost everywhere and is even;
absolute continuity proves the assertion by induction.
Its CDF is convex on $[0,W/2]$, vanishes at zero, and equals $1/2$ at $W/2$.
Thus its value at $N$ is at most $N/W$.
Since $N\le B$ and $A>B$, $N/(A-B+2N)\le B/(A+B)$.
Multiplication by $A$ proves the claim, including $B=0$.
\end{proof}

\begin{theorem}[Linear budget for the actual envelope]\label{net:envelope-upper}
For every fixed edge-owned DAG and every nonnegative real profile,
$\sum_ey_e(c)\le KT(c)$.  More precisely, if $T>0$ and
$d_C=\sum_{h\in C}D_h(c)$, then
\begin{equation}\label{net:envelope-coordinate}
 y_e(c)\le T D_e(c)\sum_{C\ni e}\frac1{d_C}.
\end{equation}
\end{theorem}
\begin{proof}
Fix $c$, with $T>0$, and a minimum cut $A$ for $D(c)$, so $d_A=T$.
Write $Z_C=\sum_{h\in C}D_h(c)U_h$ using one common uniform vector.
If $D_e=0$, all own underreports pay zero.
Reports $t\ge D_e$ leave every feasible flow unchanged, so it suffices to
consider $0<t\le D_e$.  Set $x=t/D_e$ and $c^t=(t,c_{-e})$.
Feasible-set inclusion and scaling a flow attaining each effective cap give
\[
 xD_h(c)\le D_h(c^t)\le D_h(c),\qquad D_e(c^t)=t.
\]
If the selected cut at $c^t$ is $C$, its score is no larger than that of $A$,
and hence $xZ_C\le Z_A$.  Lemma~\ref{net:uniform-comparison}, with totals
$xd_C$ and $T$, yields
\[
 t\Prob(W(c^t)=C)\le TD_e/d_C.
\]
Sum over cuts containing $e$ before taking its own supremum, obtaining
\eqref{net:envelope-coordinate}.  Summing the resulting bounds gives
$\sum_ey_e\le T\sum_C\sum_{e\in C}D_e/d_C=KT$.
For $T=0$, all effective capacities vanish.  Coverage and incentives follow
from domination and the expanding supremum domain.
\end{proof}
For comparison with earlier accounting, let $L$ be the longest path length,
$h(C)=\max_P|P\cap C|$, and $H=\sum_Ch(C)$.
Since $H\ge K$, Theorem~\ref{net:envelope-upper} immediately implies the
archived weaker bound $2HT$.  Also $H\le KL$ and $L\le K$: for each edge
of a fixed path, deleting that edge and every edge outside the path contains
a minimal cut whose trace on the path is exactly that edge, giving $L$
distinct cuts.  Hence the former $2KLT$ and $2K^2T$ conclusions are preserved.
The earlier structural estimate $H=O(K^{3/2})$ and its sharp Fibonacci-family
example concern the statistic $H$; their archived proofs are independent
of the improved envelope estimate.  They do not supply an envelope lower
bound and are not repeated here.  The sharp general mechanism theorem uses
a different payment rule.

\begin{theorem}[Growing lower bounds for the envelope]
\label{net:envelope-lower}
There are explicit fixed-owner DAG families at positive profiles with $T=1$
and growing cut count $K$ for which
\[
 \sum_ey_e=\Omega\!\left(\frac{K^{1/401}}{\ln K}\right).
\]
There are also explicit simple-DAG families with
$\sum_ey_e=\Omega(K^{1/19201}/\ln K)$.
Both statements apply to every own-report monotone pointwise majorant of $z$.
They do not apply to arbitrary truthful coalition-covering mechanisms.
\end{theorem}
\begin{proof}
First take an integer $r\ge4$ and set
\[
 N=2^r,\qquad m=400r,\qquad k=10m,\qquad
 a=6/5,\qquad b=3/10,\qquad t=b/m.
\]
Place $N$ modules in series.  Module $j$, with entry $s_j$, exit $q_j$,
vertex $A_j$ and vertices $B_{j,i}$, has edges and capacities
\[
 \begin{array}{c|c|c}
 \text{edge}&\text{endpoints}&\text{capacity}\\\hline
 e_j&s_j\to A_j&1\\
 y_{j,i}&s_j\to B_{j,i}&a/m\\
 g_{j,i}&A_j\to B_{j,i}&2\\
 h_{j,i}&B_{j,i}\to q_j&2\\
 Z_{j,\ell}&A_j\to q_j&2\quad(1\le\ell\le k).
 \end{array}
\]
After the modules add a final parallel bundle of $100000$ edges, each of
capacity $1/100000$.  All listed edges have distinct original owners.
For each $j$ separately compare the common original profile with the profile
lowering only $e_j$ to $t$.

A module has $2^{m+1}$ minimal cuts.  A cut containing $e$ chooses one of
$y_i,h_i$ in each branch.  A cut excluding $e$ contains all $Z$ edges and
chooses either $h_i$ or both $y_i,g_i$ in each branch.
These are necessary to block the direct and branch paths, sufficient to block
all paths, and minimal because each chosen edge has an unblocked-path witness
when removed.  A series-composition minimal cut lies in one component, so
\begin{equation}\label{net:aggregate-K}
 K=N2^{m+1}+1=2^{401r+1}+1.
\end{equation}
If a module's focal report is $q\in\{1,t\}$, its raw value is $a+q$:
the source cut is an upper bound, attained by routing $q$ along $e,Z$ and
each $y_i$ along $h_i$.  Thus all modules remain above one, and the final
bundle fixes global value one at every comparison.
The effective capacities inside this module are
\[
 D_e=D_{g_i}=D_{Z_\ell}=q,\quad D_{y_i}=a/m,\quad
 D_{h_i}=\min(1,q+a/m).
\]
Incoming cuts and the final bound give necessity.  Prefix/suffix routes
attain each bound; all other modules retain raw value greater than one, so
changing the focal report changes no effective capacity in those modules.

Let
\[
 A_q=\sum_{i=1}^m\min((a/m)U_{y_i},D_{h_i}U_{h_i}),
 \qquad Z_q=q\sum_{\ell=1}^kU_{Z_\ell}.
\]
The least containing score is $qU_e+A_q$; every excluding score is at least
$Z_q$.  For $0<u\le v$, integrating the product survival function gives
$\E\min(uU,vV)=u/2-u^2/(6v)$.  Therefore
\[
 \E A_{\rm high}=\frac35-\frac6{25m}\ge\frac{59}{100},
 \qquad \E A_{\rm low}=\frac{11}{25}.
\]
The summands lie in $[0,a/m]$.  The bounded-sum inequality
$\Prob(\sum(X_i-\E X_i)\ge s)\le
\exp(-2s^2/\sum_i\operatorname{range}(X_i)^2)$, and its lower-tail version,
give
\[
 \Prob(A_{\rm high}\le.53)\le e^{-m/200},\quad
 \Prob(A_{\rm low}\ge.47)\le e^{-m/800},\quad
 \Prob(Z_{\rm low}\le.53)\le e^{-9409m/4500}\le e^{-2m}.
\]
Here $\E Z_{\rm low}=3/2$ and its sum of squared ranges is $9/(10m)$;
high $Z$ stochastically dominates low $Z$.
For completeness, the bounded-sum inequality follows by convexity reducing
each variable to its endpoints: the centered Bernoulli log-mgf has second
derivative at most one quarter, so it is at most $\lambda^2/8$ after unit
scaling.  Multiply the mgf bounds and optimize the exponential Markov bound.
The final-bundle random score $Q$ has mean $1/2$ and variance $1/1200000$.
Chebyshev gives $\Prob(Q\notin[.49,.51])\le1/120$.

At the profile lowering $e_j$, union bound the event that its $A_{\rm low}$
is below $.47$, every other $A_{\rm high}$ exceeds $.53$, all $Z$ exceed
$.53$, and $Q\in[.49,.51]$.  Its failure probability is at most
\[
 e^{-m/800}+Ne^{-m/200}+Ne^{-2m}+1/120
 \le1/4+1/16+1/16+1/120=23/60.
\]
The last inequalities follow from $m=400r$, $r\ge4$, and $e>2$.
On the good event the focal containing score is at most $.47+t<.48$,
whereas its excluding and all other module scores exceed $.53$ and the
final score is at least $.49$.  The global winner therefore contains $e_j$
for every realization of its own uniform.  Thus
\[
 z_{e_j}(c^{t,j})\ge37t/60>t/2,\qquad
 \sum_ey_e(c)\ge Nt/2=\frac{3\,2^r}{8000r}.
\]
Together with \eqref{net:aggregate-K} this proves the first bound.
The separate underreport comparisons justify different own suprema at the
one original profile; they do not sum actual payments from different profiles.

For the simple construction set $m=1600r$, retaining $N,k,a,b,t$ as above.
Replace every $y_i,h_i,Z_\ell$ by a private two-edge path, each copy with
the original capacity and its own fixed owner.  Keep $e,g_i$ unsplit.
Use $2000$ private final two-edge paths in parallel, each copy of capacity
$1/2000$.  All private vertices are distinct.  There are
$N(25m+1)+4000$ owners, and the graph is simple and acyclic.
Raw values and effective-capacity bounds are unchanged for each copy.
The random capacity of a replaced path is its cap times
$V=\min(U_1,U_2)$, with mean $1/3$, variance $1/18$, and survival $(1-v)^2$.
Different paths give independent variables.  Direct integration yields
\[
 \E\min(uV_1,vV_2)=u\left[\frac13-\frac{u}{6v}
                                      +\frac{u^2}{30v^2}\right],\quad u\le v.
\]
Thus
\[
 \E A_{\rm high}=\frac25-\frac6{25m}+\frac{(6/5)^3}{30m^2}
       \ge.39,\qquad \E A_{\rm low}=166/625.
\]
The same range bounds and bounded-sum inequality give
\[
 \Prob(A_{\rm high}\le.36)\le e^{-m/800},\quad
 \Prob(A_{\rm low}\ge.30)\le e^{-1849m/1125000}\le e^{-m/800},
 \quad \Prob(Z_{\rm low}\le.36)\le e^{-1024m/1125}\le e^{-m/2}.
\]
Now $\E Z_{\rm low}=1$, with squared ranges still $9/(10m)$.
The final score has mean $1/3$ and variance $1/(18\cdot2000)$, so
$\Prob(Q\notin[47/150,53/150])\le5/72$.
The corresponding good event fails with probability at most
\[
 1/16+1/256+1/16+5/72=457/2304<1/2.
\]
Indeed $Ne^{-m/800}\le2^{-r}$, $e^{-m/800}\le2^{-2r}$, and
$Ne^{-m/2}\le2^{-r}$.  The focal score is at most
$.30+t<.31<47/150$; its excluding and all other module scores exceed
$.36>53/150$.  Hence it wins with probability greater than one half and
\[
 \sum_ey_e(c)\ge\frac{3\,2^r}{32000r}.
\]
Containing cuts now have four choices per branch.  Excluding cuts have
four per branch and two choices on each of the $k$ direct paths.  The final
paths contribute $2^{2000}$ cuts.  Therefore
\[
 K=N4^m(1+2^k)+2^{2000}
   =2^{19201r}+2^{3201r}+2^{2000},\qquad
 2^{19201r}\le K\le3\,2^{19201r},
\]
which proves the second asymptotic bound.
Finally, if a rule $p$ is own-report monotone and $p\ge z$ pointwise, then
$p_e(c)\ge p_e(t,c_{-e})\ge z_e(t,c_{-e})$ for every $t\le c_e$;
taking suprema gives $p\ge y$.
\end{proof}

\begin{remark}[Why this is a mechanism-specific separation]
\label{net:fixed-cut}
Fix a cut $C$ from topology alone and pay
$p_e=\ind\{e\in C\}\min(c_e,T(c))$.
Both arguments of the minimum are own-report nondecreasing.  For a coalition,
\[
 \sum_{e\in S}p_e\ge
 \min\left\{\sum_{e\in C\cap S}c_e,T\right\}\ge v_c(S),
\]
because either a capped summand already reaches $T$, or none is capped.
Total payment is at most $|C|T$.  The two lower-bound families above have
fixed final cuts of sizes $100000$ and $2000$, respectively, and therefore
admit constant-factor truthful covering rules on their entire report domains.
The envelope nevertheless grows without bound on them.  Its exact worst-case
growth between Theorems~\ref{net:envelope-upper} and~\ref{net:envelope-lower}
remains open.  The latter blocks the sharp-order target for every monotone
pointwise majorant of the original rule, while Theorem~\ref{prod:upper}
attains that target by departing from this class.
\end{remark}
